
\section{OA-02: BB84 basis sifting and intercept--resend}
\source{2026 Optics Assignment 2, Question 3; released solutions.}
\subsection*{Step 1: record the encoding}
Use \(Z=\{\ket H,\ket V\}\) and
\(X=\{\ket D,\ket A\}\), with
\(\ket D=(\ket H+\ket V)/\sqrt2\) and
\(\ket A=(\ket H-\ket V)/\sqrt2\).
\[
\begin{array}{c|cc}
\text{Bit}&Z\text{ basis}&X\text{ basis}\\\hline
0&\ket H&\ket D\\
1&\ket V&\ket A
\end{array}
\]
Alice and Bob choose bases independently and uniformly,
so they match with probability \(1/2\).

\subsection*{Step 2: account for transmission and detection}
For \(N\) transmitted pulses, channel transmissivity
\(\eta_{\rm ch}\) and detector efficiency \(\eta_d\),
assuming single photons and no dark counts,
\[
\boxed{\mathbb E[N_{\rm detected}]=N\eta_{\rm ch}\eta_d},
\qquad
\boxed{\mathbb E[N_{\rm sifted}]=\tfrac12
N\eta_{\rm ch}\eta_d}.
\]
For example, with \(N=10^6\), \(\eta_{\rm ch}=0.2\),
\(\eta_d=0.8\), expect 160,000 detected and
80,000 sifted bits \emph{before} sampling and key distillation.

\subsection*{Step 3: derive the intercept--resend error rate}
Condition on an Alice--Bob matched basis (the sifted key):
\[
P(\text{Eve matches Alice})=\tfrac12,\quad
P(\text{error}\mid\text{match})=0.
\]
If Eve chose the wrong conjugate basis, her result is
random; once Bob measures in Alice's basis, he has a
\(1/2\) chance of the wrong bit. Therefore
\[
\boxed{\mathrm{QBER}=\tfrac12(0)+\tfrac12(\tfrac12)=\tfrac14}.
\]
To estimate the QBER, they publicly reveal \(m\)
sifted bits. Under independent errors,
\[
P(\ge1\text{ error observed})=1-(1-Q)^m.
\]
For \(m=20,\ Q=1/4\),
\[
\boxed{P(\text{detect attack})=1-(3/4)^{20}
\simeq0.99683}.
\]
Those 20 disclosed bits must be removed from the secret
key. Error correction reconciles discrepancies; privacy
amplification reduces an eavesdropper's possible information.
\checkmethod{The \(25\%\) is calculated \emph{after basis
sifting}, not across all sent photons. The \(50\%\) sift factor
is independent of optical transmission in this ideal model.}
